% DAFN - Containers
% Roberto Masocco <roberto.masocco@uniroma2.it>
% May 1, 2022

\documentclass{beamer}

% Slides layout
\usepackage[
    title={Containers},
    subtitle={Linux Kernel and Docker},
    event={DAFN},
    author={Roberto Masocco},
    longauthor={Roberto Masocco},
    email={roberto.masocco@uniroma2.it},
    institute={Tor Vergata},
    longinstitute={University of Rome "Tor Vergata"},
    department={Department of Civil Engineering and Computer Science Engineering},
    researchgroup={Intelligent Systems Lab},
    date={May 5, 2022}
]{utvengbeamer}

% Code listings settings
\usepackage[nomath]{lmodern}
\definecolor{codegreen}{rgb}{0 0.5 0}
\definecolor{codered}{rgb}{1 0 0}
\definecolor{codeocher}{rgb}{0.8 0.47 0.13}
\usepackage{listings}
\lstdefinestyle{beamer}{
    basicstyle=\ttfamily\small,
    commentstyle=\color{codegreen},
    breakatwhitespace=false,
    captionpos=b,
    frame=lines,
    keepspaces=true,
    keywordstyle=\color{codered}\bfseries,
    numbers=left,
    numbersep=5pt,
    numberstyle=\footnotesize,
    showspaces=false,
    showstringspaces=false,
    showtabs=false,
    stringstyle=\color{codeocher},
    tabsize=2
}
\lstset{style=beamer}
\lstdefinelanguage{Dockerfile}{
  alsoletter={[, ], _, /},
  morecomment=[l][\color{codegreen}]{\#},
  morekeywords={FROM, RUN, ADD, COPY, LABEL, ENV, ARG, CMD}
}
\lstdefinelanguage{compose}{
  alsoletter={:, -, /},
  morecomment=[l][\color{codegreen}]{\#},
  morekeywords={services:, build:, context:, network_mode:, args:, environment:, command:, volumes:, -}
}

\usepackage{hyperref}

\begin{document}

% --- Title page ---
\frame{\titlepage}

% --- Disclaimer ---
\begin{frame}{Disclaimer}
\begin{block}{}
  \centering
  What follows is heavily based on specific features of the Linux kernel.\\
  \textbr{Compatibility with different platforms cannot be guaranteed.}
\end{block}
\end{frame}

% --- Table of contents ---
\begin{frame}
\frametitle{Roadmap}
\tableofcontents
\end{frame}

% --- Section 1 ---
% Section 1 - Containers
% Roberto Masocco <roberto.masocco@uniroma2.it>
% May 1, 2022

% ### Containers ###
\section{Containers}
\graphicspath{{figs/section1/}}

% --- Why containers? ---
\begin{frame}{Why containers?}
\begin{exampleblock}{Example: Packaging Applications}
  Suppose you are ready to distribute your new application:
  \begin{itemize}
    \item you need to be sure that it is compatible with all the \textbf{platforms} you chose to support;
    \item you need to figure out a way to deal with \textbf{dependencies};
    \item you want to publish some kind of self-contained, easily-identifiable \textbf{package}.
  \end{itemize}
\end{exampleblock}
\end{frame}
\begin{frame}{Why containers?}
\begin{exampleblock}{Example: Isolating Applications}
  Suppose you are deploying applications on a server:
  \begin{itemize}
    \item you want to define \textbf{resource quotas} and \textbf{permissions} for each;
    \item you want to be sure that each module has what it needs to operate, but \textbf{nothing more};
    \item you want to \textbf{isolate} each module for security reasons, in case something goes wrong.
  \end{itemize}
\end{exampleblock}
\end{frame}
\begin{frame}{Why containers?}
\begin{exampleblock}{Example: Replicating Environments}
  Suppose you are developing applications for a specific system (maybe with a different architecture):
  \begin{itemize}
    \item you want to have a \textbf{local copy} of such system without carrying one with you;
    \item you want to have all \textbf{libraries} and \textbf{dependencies} installed without tainting your own system;
    \item you would like to \textbf{deploy} the entire installation with just a few commands, without running any script but simply copying data.
  \end{itemize}
\end{exampleblock}
\end{frame}
\begin{frame}{Why containers?}
\begin{columns}
  \column{.5\textwidth}
  A possible solution to many of the previous situations could be a set of \textbg{virtual machines}.\\
  However, virtual machines are \textbg{slow}, hypervisors take up \textbg{system resources} and guest kernels must always be \textbg{tweaked}.\\
  In each of the above scenarios something simpler would be enough, especially since \textbg{the OS is not involved}, only applications are.\\
  This is what a \textbg{container} is.

  \column{.5\textwidth}
  \begin{figure}
    \centering
    \includegraphics[scale=.7]{freebsdjail.png}
    \label{fig:jail}
    \caption{FreeBSD jail logo}
  \end{figure}
\end{columns}
\end{frame}

% --- Containers and the Linux kernel ---
\begin{frame}{Containers and the Linux kernel}
\end{frame}


% --- Section 2 ---
% Section 2 - Docker
% Roberto Masocco <roberto.masocco@uniroma2.it>
% May 1, 2022

% ### Docker ###
\section{Docker}
\graphicspath{{figs/section2/}}


% --- Sectioin 3 ---
% Section 3 - Docker Compose
% Roberto Masocco <roberto.masocco@uniroma2.it>
% May 1, 2022

% ### Docker Compose ###
\section{Docker Compose}
\graphicspath{{figs/section3/}}


\end{document}
